At precisely these original quartet degrees, the complete
ODD1--14 calculation retains the degree in the finite DMR error.
Write $q=q_k=(k+1)^2[1+k(m-1)]$ in the next displays and put
\[
\begin{gathered}
 \eta=(p-1)/p,\quad
 C_r=1+\frac{\alpha+\log(M_\alpha/\vartheta_h)+B}{\log2},\quad
 C_d=2C_r+M/3,\quad C_L=1+C_d,\\
 C_E=6C_r^2+(2M+4)C_d/3,\quad
 C_s=\left|\log\frac{C_h}{c_\Gamma b_h}+2\log M_\alpha\right|,
 \quad C_{\rm tail}=\frac{2C_h}{(p-1)M_1},\\
 M_1=\int_{-1}^1e^{\alpha y}w_h(y)\,dy,\quad
 T_k^{\rm deg}=\log(4m(C_L+2))+3\log(k+2),\\
 z_k=\frac{kT_k^{\rm deg}}q,\qquad
 R_k^*=\max(1,z_k^{-1/p}),\qquad
 g_0=\log\sqrt{2\pi}.
\end{gathered}\tag{HMR44a}
\]
Here $w_h=w$ is the original source density and $M_1$ is its
tilted finite mass, whereas $\vartheta_h$ is its un-tilted
finite mass. The scalar $T_k^{\rm deg}$ is exactly ODD3's
$T_k$; its superscript prevents confusion with the original
HC threshold $4q\log(D_hk)$. No source variable is changed.
Every displayed constant except $T_k^{\rm deg},z_k,R_k^*$
is independent of $k$ and of the cutoff radius.

For $N\in\{q-1,q,2q-1,2q\}$ define the explicit finite
degree-dependent majorant and its common bound by
\[
\begin{aligned}
 \Psi_{k,N}^{\rm deg}={}&\frac{N+1}{q}
 \left\{C_{\rm tail}\min(2^{1-p},z_k^\eta)
       +\frac{(p+1)\log k+C_s}{k}\right\}
       +\frac{2MN T_k^{\rm deg}}{kq}\\
 &+C_Ez_k+C_d\max(z_k,z_k^\eta)
                     +\frac{(g_0+2)C_d}{q},\\
 \Psi_k={}&\left(2+\frac1q\right)
 \left\{C_{\rm tail}\min(2^{1-p},z_k^\eta)
       +\frac{(p+1)\log k+C_s+2MT_k^{\rm deg}}{k}\right\}\\
 &+C_Ez_k+C_d\max(z_k,z_k^\eta)
                     +\frac{(g_0+2)C_d}{q},\\
 0\leq{}&\frac{\Xi_{k,N}}{kq}
        \leq\Psi_{k,N}^{\rm deg}\leq\Psi_k,\\
 D_k^-={}&kq(\Psi_{k,q-1}^{\rm deg}+\Psi_{k,q}^{\rm deg}),\qquad
 D_k^+=kq(\Psi_{k,2q-1}^{\rm deg}+\Psi_{k,2q}^{\rm deg}),\\
 \delta_k^\sigma={}&F_{\mathrm{ar},k}-F_{\sigma,k},\qquad
 \delta_k^\sigma,\Delta F_K,\Delta F_B\in[-D_k^-,D_k^+],\\
 \max\{|\delta_k^\sigma|,|\Delta F_K|,|\Delta F_B|\}
 &\leq2kq\Psi_k,\\
 \frac{\Xi_{k,N}}{kq},\quad
 \frac{\max\{|\delta_k^\sigma|,|\Delta F_K|,|\Delta F_B|\}}{kq}
 & =O_h\!\left(\frac{\log k}{k}
                 +\left(\frac{k\log k}{q}\right)^\eta\right).
\end{aligned}\tag{HMR44}
\]
All finite inequalities hold for $k\geq3$ and every $m\geq1$.
Here is their proof, including the retained degree factors.
The complete DMR29 and DMR32 bounds inserted into HMR38 give,
for every $R\geq1$, exactly ODD7:
\[
\begin{aligned}
 \frac{\Xi_{k,N}}{kq}\leq{}&\frac{N+1}{q}
 \left\{C_{\rm tail}(1+R)^{1-p}
                 +\frac{(p+1)\log k+C_s}{k}\right\}
            +\frac{2MN T_k^{\rm deg}}{kq}\\
 &+\frac{kT_k^{\rm deg}}q(C_E+C_dR)
                     +\frac{(g_0+2)C_d}{q}.
\end{aligned}
\]
The original tilted-tail estimate is
$\log(M_\alpha/M_R)\leq C_{\rm tail}(1+R)^{1-p}$.
The error bound uses the actual $r\leq C_rk$,
$d\leq C_dk$, $Y\leq kR$, $L\leq C_LN$ and
$\log[2(L+2)]\leq T_k^{\rm deg}$ on these quartet degrees.
Thus it is the same full divided-jet determinant calculation,
including its central rational multiplier, rather than a
replacement source or a scalar-moment comparison.
If $z_k\leq1$ then $z_kR_k^*=z_k^\eta$; if $z_k\geq1$
then $z_kR_k^*=z_k$. In both branches
$(1+R_k^*)^{1-p}\leq\min(2^{1-p},z_k^\eta)$.
Substitution proves the bound by $\Psi_{k,N}^{\rm deg}$.
The inequalities $N/q\leq(N+1)/q\leq2+1/q$ prove its
bound by the complete ODD5 majorant $\Psi_k$.
Now use the two negative and two positive terms in HMR43
separately. This proves $[-D_k^-,D_k^+]$ for the original
quotient and each of its original kernel and boundary factors.
The stronger exact interval in HMR43 is retained; each finite
majorant here bounds its corresponding actual endpoint sum.
The common bound $2kq\Psi_k$ follows because each side has
two terms. Since $q\geq(k+1)^2$, $z_k=O_h(k\log k/q)\to0$;
then $z_k\leq z_k^\eta$ and $q^{-1}=O(k^{-2})$, proving
the stated rate directly from the finite formulas.

There is a sharper signed consequence at the original
multiple-zero degree. Put
$C_m=p+1+6M=32m+245/2$. If $m\geq2$ and the actual
endpoint sequence satisfies $N/q\to s\in\{1,2\}$, then
\[
\begin{gathered}
 \limsup_{k\to\infty}\frac{\Xi_{k,N}}{q\log k}
       \leq sC_m,\qquad
 \lim_{k\to\infty}\frac{D_k^-}{q\log k}=2C_m=64m+245,\\
 \lim_{k\to\infty}\frac{D_k^+}{q\log k}=4C_m=128m+490,\\
 -(64m+245)\leq\liminf_{k\to\infty}\frac{X_k}{q\log k}
 \leq\limsup_{k\to\infty}\frac{X_k}{q\log k}
 \leq128m+490,\qquad
 X_k\in\{\delta_k^\sigma,\Delta F_K,\Delta F_B\}.
\end{gathered}\tag{HMR44b}
\]
To prove these claims, retain $q=(m-1)k^3+O_m(k^2)$ and
$p\geq23/2$, so $\eta>1/2$.
After multiplication of $\Psi_{k,N}^{\rm deg}$ by
$k/\log k$, its tail and cutoff terms tend to zero because
$kz_k^\eta/\log k
=O_m(k^{1-2\eta}(\log k)^{\eta-1})\to0$,
$kz_k/\log k=O_m(k^{-1})$, and
$k/(q\log k)\to0$. Its two remaining terms tend to
$s(p+1)$ and $6Ms$, since $T_k^{\rm deg}/\log k\to3$.
Therefore $k\Psi_{k,N}^{\rm deg}/\log k\to sC_m$.
Summing the two low and the two high endpoints proves
the displayed limits for $D_k^-,D_k^+$ and the signed
bounds by HMR44. For $m=1$ retain $q=(k+1)^2$ and
$\eta=5/7$; HMR44 instead gives the full rate
$O_h((\log k/k)^{5/7})$ on the $kq$ scale.
No $q\log k$ bound in HMR44b is assigned to that branch.
Nor is this quartet-degree estimate assigned to the
simple-pair degree $N=2k+2$; the all-degree HMR38 remains
available there with its literal finite error.

The original HC5 threshold and its nonnegative residual
also receive these same signed endpoints, with their
original constant $D_h>0$:
\[
\begin{gathered}
 \mathcal R_k=F_{\sigma,k}+\delta_k^\sigma
                         -4q\log(D_hk)\geq0,\\
 \max\{0,F_{\sigma,k}-4q\log(D_hk)-D_k^-\}
 \leq\mathcal R_k
 \leq F_{\sigma,k}-4q\log(D_hk)+D_k^+.
\end{gathered}\tag{HMR44c}
\]
Substitution of HMR44 into the first exact identity proves
both endpoints. This retains the full fixed-Gamma functional
$F_{\sigma,k}$ and the complete HC residual; the independent
nonnegativity of that residual gives precisely the lower maximum.

